\renewcommand*{\authorOfNextLines}{Helmut Quinz}
Direkte Zitate sind aus der Quelle wörtlich übernommenes Textfragmente.
Je nach Länge des Zitats wird zwischen
\begin{itemize}
	\item Kurzzitat und
	\item Blockzitat unterschieden.
\end{itemize} 
Für das direkte Zitat sind folgende Regeln zu berücksichtigen:
\begin{itemize}
	\item Die Zitate sind \textbf{buchstabengetreu} zu anzuführen. 
		Selbst allfällige Fehler des Ursprungsdokuments sind zu übernehmen, dürfen aber mit \gqq{sic} gekennzeichnet werden.
	\item Die Zitate sind in Anführungszeichen zu setzen. 
		Bei eingerückten Blockkommentaren dürfen diese im Zuge Ihrer Diplomarbeit entfallen.
	\item Um direkte Zitate besser vom umgebenden Text unterscheiden zu können, wird darum gebeten diese Kursiv zu setzen (\verb|\textit{...}|)
	\item Beim direkten zitieren aus Büchern ist die Seitenzahl anzugeben
	\item Direkte Zitate mit mehr als 40 Wörtern sind als eingerücktes Blockzitat auszuführen.
		Unter \LaTeX ~eignet sich dazu die \texttt{quote}-Umgebung:\newline
		 (\verb|\begin{quote} ... \end{quote}|)
	\item Die Quellenangabe steht bei Kurzzitaten nach dem Anführungszeichen, jedoch vor dem Satzzeichen, bei Blockzitaten nach dem Satzzeichen.
			
\end{itemize}

\subsubsection{Beispiel für ein Kurzzitat}
Das Nehmen des Lebens als Bestrafung wird auch in der Literatur reichlich behandelt.
So lässt J.R.R Tolikien in seinem Werk \gqq{Herr der Ringe} Gandalf sagen: 
\gqqEmph{Viele, die leben, verdienen den Tod. Und manche, die sterben, verdienen das Leben. Kannst du es ihnen geben? Dann sei auch nicht so rasch mit einem Todesurteil bei der Hand.} (\cite[Seite 74]{Tolkien1954})
\subsubsection{Beispiel für ein Blockzitat}
Die Bedeutung der Zahl \gqq{\texttt{42}} in der Informationstechnologie durch folgenden Abschnitt aus einem Meisterwerk der Weltliteratur verdeutlicht:
\begin{quote}
	\textit{Beide Männer waren auf diesen Augenblick gedrillt worden, 
			ihr Leben war eine einzige Vorbereitung auf diesen Moment gewesen, 
			man hatte sie bereits bei ihrer Geburt als diejenigen ausgewählt, 
			die der Antwort beiwohnen würden, aber selbst sie wurden gewahr, dass sie jetzt nach Luft schnappten und rumhampelten wie aufgeregte Kinder.\\
			>>Und du bist bereit, sie uns zu geben?<<, drängte Luunquoal.\\
			>>Das bin ich.<<\\
			>>Jetzt?<<\\
			>>Jetzt<<, sagte Deep Thought.\\
			Beide Männer leckten sich ihre trockenen Lippen.\\
			>>Allerdings glaube ich nicht<<, setzte Deep Thought hinzu, >>dass sie euch gefallen wird.<<\\
			>>Das macht doch nichts!<<, sagte Phouchg. >>Wir müssen sie nur jetzt erfahren. Jetzt!<<\\
			>>Jetzt?<<, fragte Deep Thought.\\
			>>Ja! Jetzt …<<\\
			>>Also schön<<, sagte der Computer und versank wieder in Schweigen. Die beiden Männer zappelten nervös hin und her. Die Spannung war unerträglich.\\
			>>Sie wird euch bestimmt nicht gefallen<<, bemerkte Deep Thought.\\
			>>Sag sie uns trotzdem!<<\\
			>>Na schön<<, sagte Deep Thought. >>Die Antwort auf die Große Frage \dots<<\\
			>>Ja \dots !<<\\
			>>… nach dem Leben, dem Universum und allem \dots<<, sagte Deep Thought.\\
			>>Ja \dots !<<\\
			>>\dots ~lautet \dots <<, sagte Deep Thought und machte eine Pause.\\
			>>Ja \dots !<<\\
			>>\dots ~lautet \dots <<\\
			>>Ja \dots !\dots ???<<\\
			>>Zweiundvierzig<<, sagte Deep Thought mit unsagbarer Erhabenheit und Ruhe.} \cite[Seiten 186 ff]{Douglas2013}
\end{quote}